%%%%%%%%%%%%%%%%%%%%%%%%%%%%%%%%%%%%%%%%%%%%%%%%%%%%%%%%%%%%%%%%%%%
% Standalone university LaTeX document -- Physik IV Blatt 03
%%%%%%%%%%%%%%%%%%%%%%%%%%%%%%%%%%%%%%%%%%%%%%%%%%%%%%%%%%%%%%%%%%%
\documentclass[11pt,a4paper]{scrartcl}

%%%%%%%%%%%%%%%%%%%%%%%%%%%%%%%%%%%%%%%%%%%%%%%%%%%%%%%%%%%%%%%%%%%
% Embedded file: includes/university_metadata.tex
%%%%%%%%%%%%%%%%%%%%%%%%%%%%%%%%%%%%%%%%%%%%%%%%%%%%%%%%%%%%%%%%%%%
\newcommand{\CourseTitle}{Proseminar zur Physik IV: Kerne und Teilchen}
\newcommand{\CourseShortTitle}{Physik IV}
\newcommand{\DocumentKind}{Hausaufgabe, Tafelvortrag und Lernlösung}
\newcommand{\DocumentTitle}{Aufgabenblatt 3}
\newcommand{\DocumentSubtitle}{Teilchenbeschleuniger und Luminosität des LHC}
\newcommand{\AuthorName}{Tobias Laser}
\newcommand{\InstitutionName}{Universität Innsbruck}
\newcommand{\SubmissionDate}{24. März 2026}

\newcommand{\makeuniversitytitle}{%
    \begin{titlepage}
        \centering
        \vspace*{2cm}
        {\Large \CourseTitle\par}
        \vspace{1.2cm}
        {\huge\bfseries \DocumentTitle\par}
        \ifx\DocumentSubtitle\empty\else
            \vspace{0.4cm}
            {\Large \DocumentSubtitle\par}
        \fi
        \vspace{1cm}
        {\Large \DocumentKind\par}
        \vfill
        {\large \AuthorName\par}
        \vspace{0.2cm}
        {\large \InstitutionName\par}
        \vspace{0.8cm}
        {\large \SubmissionDate\par}
    \end{titlepage}
    \pagenumbering{arabic}
}

\newcommand{\makechaptertitle}{%
    \begin{center}
        {\Large \CourseTitle\par}
        \vspace{0.5cm}
        {\huge\bfseries \DocumentTitle\par}
        \ifx\DocumentSubtitle\empty\else
            \vspace{0.25cm}
            {\Large \DocumentSubtitle\par}
        \fi
        \vspace{0.5cm}
        {\large \AuthorName \hfill \SubmissionDate\par}
    \end{center}
    \vspace{0.5cm}
}
\newcommand{\makephysiktitle}{\makeuniversitytitle}

%%%%%%%%%%%%%%%%%%%%%%%%%%%%%%%%%%%%%%%%%%%%%%%%%%%%%%%%%%%%%%%%%%%
% Embedded file: includes/university_packages.tex
%%%%%%%%%%%%%%%%%%%%%%%%%%%%%%%%%%%%%%%%%%%%%%%%%%%%%%%%%%%%%%%%%%%
\usepackage[margin=2.5cm]{geometry}
\usepackage[onehalfspacing]{setspace}
\usepackage{scrlayer-scrpage}

\usepackage[T1]{fontenc}
\usepackage[utf8]{inputenc}
\usepackage[ngerman]{babel}
\usepackage{lmodern}
\usepackage{microtype}
\usepackage{csquotes}

\usepackage{amsmath,amsthm,amssymb,mathtools}
\usepackage{bm}
\usepackage{icomma}
\usepackage{cancel}

\usepackage{graphicx}
\usepackage{tikz}
\usetikzlibrary{arrows.meta,calc,decorations.pathmorphing,positioning,patterns,angles,quotes}
\usepackage{pgfplots}
\pgfplotsset{compat=1.18}

\usepackage[locale=DE,space-before-unit=true,per-mode=symbol,separate-uncertainty=true]{siunitx}
\usepackage{booktabs,multirow,array,tabularx,longtable}
\usepackage{caption,subcaption}
\usepackage{placeins}
\usepackage{enumitem}
\usepackage{xparse}

\usepackage[most]{tcolorbox}
\tcbuselibrary{breakable,skins,theorems}

\usepackage[breaklinks=true,colorlinks=true,linkcolor=blue,urlcolor=blue,citecolor=blue]{hyperref}
\usepackage[nameinlink,noabbrev]{cleveref}

\clearpairofpagestyles
\ihead{\CourseShortTitle}
\ohead{\DocumentKind}
\cfoot{\pagemark}
\pagestyle{scrheadings}

\setlist[enumerate]{itemsep=0.4em,topsep=0.4em}
\setlist[itemize]{itemsep=0.25em,topsep=0.35em}
\captionsetup{font=small,labelfont=bf}

%%%%%%%%%%%%%%%%%%%%%%%%%%%%%%%%%%%%%%%%%%%%%%%%%%%%%%%%%%%%%%%%%%%
% Embedded file: includes/university_macros.tex
%%%%%%%%%%%%%%%%%%%%%%%%%%%%%%%%%%%%%%%%%%%%%%%%%%%%%%%%%%%%%%%%%%%
\newcommand{\R}{\mathbb{R}}
\newcommand{\N}{\mathbb{N}}
\newcommand{\Z}{\mathbb{Z}}
\newcommand{\Q}{\mathbb{Q}}
\newcommand{\C}{\mathbb{C}}

\newcommand{\set}[1]{\left\{ #1 \right\}}
\newcommand{\abs}[1]{\left| #1 \right|}
\newcommand{\norm}[1]{\left\lVert #1 \right\rVert}
\newcommand{\avg}[1]{\left\langle #1 \right\rangle}
\newcommand{\paren}[1]{\left( #1 \right)}
\newcommand{\bracket}[1]{\left[ #1 \right]}
\newcommand{\ol}[1]{\overline{#1}}

\newcommand{\dd}{\,\mathrm{d}}
\newcommand{\dv}[2]{\frac{\dd #1}{\dd #2}}
\newcommand{\pdv}[2]{\frac{\partial #1}{\partial #2}}
\newcommand{\pddv}[2]{\frac{\partial^2 #1}{\partial #2^2}}
\newcommand{\evalat}[2]{\left. #1 \right|_{#2}}

\newcommand{\vect}[1]{\bm{#1}}
\newcommand{\unitvect}[1]{\hat{\bm{#1}}}
\newcommand{\grad}{\vect{\nabla}}
\newcommand{\divergence}{\grad\!\cdot}
\newcommand{\curl}{\grad\!\times}

\newcommand{\half}{\frac{1}{2}}
\newcommand{\third}{\frac{1}{3}}
\newcommand{\twothirds}{\frac{2}{3}}
\newcommand{\hbarc}{\hbar c}
\newcommand{\alphaf}{\alpha}
\newcommand{\eV}{\si{eV}}
\newcommand{\MeV}{\si{MeV}}
\newcommand{\GeV}{\si{GeV}}
\newcommand{\TeV}{\si{TeV}}

\newcommand{\nuclide}[3]{^{#1}_{#2}\mathrm{#3}}
\newcommand{\isotope}[2]{^{#1}\mathrm{#2}}
\newcommand{\anti}[1]{\overline{#1}}
\newcommand{\pbar}{\anti{p}}
\newcommand{\nbar}{\anti{n}}
\newcommand{\ee}{e^-}
\newcommand{\ep}{e^+}
\newcommand{\mup}{\mu^+}
\newcommand{\mum}{\mu^-}
\newcommand{\nue}{\nu_e}
\newcommand{\numu}{\nu_\mu}
\newcommand{\nutau}{\nu_\tau}
\newcommand{\pionp}{\pi^+}
\newcommand{\pionm}{\pi^-}
\newcommand{\pionz}{\pi^0}
\newcommand{\kaonp}{K^+}
\newcommand{\kaonm}{K^-}
\newcommand{\kaonz}{K^0}

\newcommand{\diffxs}{\frac{\dd\sigma}{\dd\Omega}}
\newcommand{\totalxs}{\sigma_{\mathrm{total}}}
\newcommand{\lum}{\mathcal{L}}
\newcommand{\smand}{s}
\newcommand{\tmand}{t}
\newcommand{\umand}{u}
\newcommand{\Ecm}{E_{\mathrm{CM}}}

\newcommand{\ie}{d.\,h.}
\newcommand{\eg}{z.\,B.}
\newcommand{\wrt}{bezüglich}

\DeclareSIUnit{\barn}{b}
\DeclareSIUnit{\millibarn}{mb}
\DeclareSIUnit{\microbarn}{\micro b}
\DeclareSIUnit{\nanobarn}{nb}
\DeclareSIUnit{\fermi}{fm}
\DeclareSIUnit{\year}{a}
\DeclareSIUnit{\femtobarn}{fb}

%%%%%%%%%%%%%%%%%%%%%%%%%%%%%%%%%%%%%%%%%%%%%%%%%%%%%%%%%%%%%%%%%%%
% Embedded file: includes/university_boxes.tex
%%%%%%%%%%%%%%%%%%%%%%%%%%%%%%%%%%%%%%%%%%%%%%%%%%%%%%%%%%%%%%%%%%%
\tcbset{
    unibox/.style={
        enhanced,
        breakable,
        sharp corners,
        boxrule=0.6pt,
        left=1.2mm,
        right=1.2mm,
        top=1mm,
        bottom=1mm,
        before skip=0.8em,
        after skip=0.8em,
        fonttitle=\bfseries
    }
}

\newtcolorbox{calculationbox}[1][]{unibox,title=Rechnung,colback=black!2,colframe=black!55,#1}
\newtcolorbox{sidecalcbox}[1][]{unibox,title=Nebenrechnung,colback=black!1,colframe=black!40,#1}
\newtcolorbox{solutionbox}[1][]{unibox,title=Lösung,colback=blue!2,colframe=blue!45!black,#1}
\newtcolorbox{answerbox}[1][]{unibox,title=Antwort,colback=green!3,colframe=green!45!black,#1}
\newtcolorbox{notebox}[1][]{unibox,title=Notiz,colback=yellow!6,colframe=yellow!45!black,#1}
\newtcolorbox{definitionbox}[1][]{unibox,title=Definition,colback=cyan!3,colframe=cyan!45!black,#1}
\newtcolorbox{warningbox}[1][]{unibox,title=Achtung,colback=red!3,colframe=red!55!black,#1}
\newtcolorbox{summarybox}[1][]{unibox,title=Zusammenfassung,colback=violet!3,colframe=violet!50!black,#1}
\newtcolorbox{formulabox}[1][]{unibox,title=Formel,colback=orange!4,colframe=orange!65!black,#1}
\newtcolorbox{taskbox}[1][]{unibox,title=Aufgabe,colback=white,colframe=black!75,#1}
\newtcolorbox{presentationbox}[1][]{unibox,title=Tafelvortrag,colback=teal!3,colframe=teal!55!black,#1}
\newtcolorbox{learningbox}[1][]{unibox,title=Lernen,colback=purple!2,colframe=purple!50!black,#1}

%%%%%%%%%%%%%%%%%%%%%%%%%%%%%%%%%%%%%%%%%%%%%%%%%%%%%%%%%%%%%%%%%%%
% Embedded file: includes/university_theorems.tex
%%%%%%%%%%%%%%%%%%%%%%%%%%%%%%%%%%%%%%%%%%%%%%%%%%%%%%%%%%%%%%%%%%%
\theoremstyle{definition}
\newtheorem{definition}{Definition}[section]
\newtheorem{example}{Beispiel}[section]
\theoremstyle{plain}
\newtheorem{satz}{Satz}[section]
\newtheorem{lemma}{Lemma}[section]
\theoremstyle{remark}
\newtheorem*{remark}{Bemerkung}

%%%%%%%%%%%%%%%%%%%%%%%%%%%%%%%%%%%%%%%%%%%%%%%%%%%%%%%%%%%%%%%%%%%
\begin{document}

\makeuniversitytitle
\tableofcontents
\newpage

%%%%%%%%%%%%%%%%%%%%%%%%%%%%%%%%%%%%%%%%%%%%%%%%%%%%%%%%%%%%%%%%%%%
\section*{Vorwort zur Verwendung}
\addcontentsline{toc}{section}{Vorwort zur Verwendung}

Dieses Dokument ist so aufgebaut, dass es gleichzeitig als saubere Abgabe, als Vorbereitung für einen kurzen Tafelvortrag und als ausführliche Lernlösung funktioniert. Die Rechnungen sind bewusst etwas ausführlicher als in einer Minimalabgabe, damit die physikalischen Zusammenhänge sichtbar bleiben.

\begin{warningbox}[title=Wichtige Korrektur gegenüber der ursprünglichen Abgabe]
Bei Aufgabe 1(a) ist die Energie nach \num{100000} Beschleunigungsschritten nicht \(\SI{25}{\MeV}\), sondern
\(\SI{25}{\GeV}\), denn
\[
    10^5 \cdot \SI{250}{\kilo\electronvolt}=\SI{25e9}{\electronvolt}=\SI{25}{\GeV}.
\]
Die Aussage zur halben Lichtgeschwindigkeit bleibt aber inhaltlich gleich: Nach einem einzigen Beschleunigungsschritt ist die notwendige Energie bereits überschritten.
\end{warningbox}

%%%%%%%%%%%%%%%%%%%%%%%%%%%%%%%%%%%%%%%%%%%%%%%%%%%%%%%%%%%%%%%%%%%
\section{Aufgabe 1: Teilchenbeschleuniger}

\begin{taskbox}
Ein Linearbeschleuniger am SLAC arbeitet mit einer Beschleunigungsspannung von \(U=\SI{250}{\kilo\volt}\) und etwa \(N=\num{100000}\) Driftröhren. Danach werden Driftlängen für Protonen, die Ablenkung in einem Dipolmagneten und ein hypothetisches Mond-Synchrotron betrachtet.
\end{taskbox}

\subsection{Teil (a): Energiegewinn im Linearbeschleuniger}

\begin{definitionbox}[title=Prinzip]
Ein Teilchen mit Ladung \(q\), das eine Beschleunigungsspannung \(U\) durchläuft, gewinnt die kinetische Energie
\[
    \Delta E_{\mathrm{kin}} = qU.
\]
Für Elektronen ist \(q=e\). Wird der Beschleunigungsvorgang \(N\)-mal wiederholt, gilt idealisiert
\[
    E_{\mathrm{kin}} = N e U.
\]
In Elektronenvolt ist die Rechnung besonders einfach: Ein Elektron gewinnt bei \(\SI{250}{\kilo\volt}\) genau \(\SI{250}{\kilo\electronvolt}\).
\end{definitionbox}

\begin{calculationbox}[title=Energie nach \(10^5\) Beschleunigungsschritten]
\[
    E_{\mathrm{kin}}
    = N e U
    = 10^5 \cdot \SI{250}{\kilo\electronvolt}
    = \SI{25e6}{\kilo\electronvolt}
    = \SI{25e9}{\electronvolt}
    = \SI{25}{\GeV}.
\]
\end{calculationbox}

\begin{answerbox}
Der Linearbeschleuniger konnte idealisiert eine Elektronenenergie von
\[
    \boxed{E_{\mathrm{kin}} \approx \SI{25}{\GeV}}
\]
erreichen.
\end{answerbox}

\subsection{Teil (a): Anzahl der Driftröhrendurchgänge bis \(v\approx c/2\)}

\begin{learningbox}
Für \(v=0{,}5c\) ist eine relativistische Rechnung sinnvoll, weil die klassische Formel \(E_{\mathrm{kin}}=\frac12 mv^2\) bereits merklich abweicht. Relativistisch gilt
\[
    E_{\mathrm{kin}}=(\gamma-1)m_ec^2,
    \qquad
    \gamma = \frac{1}{\sqrt{1-v^2/c^2}}.
\]
Die Ruheenergie des Elektrons ist
\[
    m_ec^2 \approx \SI{511}{\kilo\electronvolt}.
\]
\end{learningbox}

\begin{calculationbox}
Für \(v=0{,}5c\) ist
\[
    \gamma = \frac{1}{\sqrt{1-(0{,}5)^2}}
    = \frac{1}{\sqrt{0{,}75}}
    \approx 1{,}1547.
\]
Damit braucht ein Elektron die kinetische Energie
\[
    E_{\mathrm{kin}}
    = (1{,}1547-1)\,\SI{511}{\kilo\electronvolt}
    \approx \SI{79,1}{\kilo\electronvolt}.
\]
Ein Beschleunigungsschritt liefert aber bereits
\[
    eU = \SI{250}{\kilo\electronvolt}.
\]
Formal wäre die benötigte Schrittzahl
\[
    n = \frac{\SI{79,1}{\kilo\electronvolt}}{\SI{250}{\kilo\electronvolt}}
    \approx 0{,}316.
\]
Da es nur ganze Durchgänge gibt, ist nach dem ersten Durchgang die halbe Lichtgeschwindigkeit bereits überschritten.
\end{calculationbox}

\begin{answerbox}
Die halbe Lichtgeschwindigkeit wird bereits nach
\[
    \boxed{1\text{ Beschleunigungsschritt}}
\]
erreicht beziehungsweise überschritten.
\end{answerbox}

\subsection{Teil (b): Driftröhrenlängen für Protonen}

\begin{definitionbox}[title=Warum Driftröhren länger werden]
Im HF-Linearbeschleuniger müssen die Teilchen während der falschen Phase des elektrischen Feldes abgeschirmt sein. Eine Driftröhre wird daher so lang gewählt, dass das Teilchen ungefähr eine halbe HF-Periode in ihr verbringt:
\[
    l = v\frac{T}{2}=\frac{v}{2f}=\frac{\beta c}{2f}.
\]
Mit der kinetischen Energie \(E_{\mathrm{kin}}\) erhält man
\[
    \gamma = 1 + \frac{E_{\mathrm{kin}}}{m_pc^2},
    \qquad
    \beta = \sqrt{1-\frac{1}{\gamma^2}},
\]
wobei
\[
    m_pc^2 \approx \SI{938,272}{\MeV}.
\]
\end{definitionbox}

\subsubsection*{(i) Protonenenergie \(\SI{9}{\MeV}\)}
\addcontentsline{toc}{subsubsection}{(i) Protonenenergie 9 MeV}

\begin{calculationbox}
\[
    \gamma = 1+\frac{\SI{9}{\MeV}}{\SI{938,272}{\MeV}}
    \approx 1{,}00959.
\]
\[
    \beta = \sqrt{1-\frac{1}{\gamma^2}}
    \approx 0{,}1375.
\]
Mit \(f=\SI{222}{\mega\hertz}\) folgt
\[
    l = \frac{0{,}1375\,c}{2\cdot \SI{222}{\mega\hertz}}
    \approx \SI{0,0929}{\metre}.
\]
\end{calculationbox}

\begin{answerbox}
Bei \(\SI{9}{\MeV}\) müssen die Driftröhren ungefähr
\[
    \boxed{l\approx \SI{9,3}{\centi\metre}}
\]
lang sein.
\end{answerbox}

\subsubsection*{(ii) Protonenenergie \(\SI{9}{\GeV}\)}
\addcontentsline{toc}{subsubsection}{(ii) Protonenenergie 9 GeV}

\begin{calculationbox}
Zuerst wird \(\SI{9}{\GeV}=\SI{9000}{\MeV}\) eingesetzt:
\[
    \gamma = 1+\frac{\SI{9000}{\MeV}}{\SI{938,272}{\MeV}}
    \approx 10{,}592.
\]
\[
    \beta = \sqrt{1-\frac{1}{\gamma^2}}
    \approx 0{,}9955.
\]
Damit ist
\[
    l = \frac{0{,}9955\,c}{2\cdot \SI{222}{\mega\hertz}}
    \approx \SI{0,672}{\metre}.
\]
\end{calculationbox}

\begin{answerbox}
Bei \(\SI{9}{\GeV}\) müssen die Driftröhren ungefähr
\[
    \boxed{l\approx \SI{67,2}{\centi\metre}}
\]
lang sein.
\end{answerbox}

\begin{notebox}[title=Physikalische Interpretation]
Bei niedriger Energie sind Protonen noch relativ langsam, daher sind die Driftröhren kurz. Bei \(\SI{9}{\GeV}\) bewegen sie sich fast mit Lichtgeschwindigkeit; die Driftröhrenlänge nähert sich deshalb dem Grenzwert \(c/(2f)\approx\SI{67,5}{\centi\metre}\).
\end{notebox}

\subsection{Teil (c): Ablenkung eines \(\SI{6,9}{\TeV}\)-Protonenstrahls}

\begin{formulabox}
Für ein relativistisches Teilchen im Dipolmagneten gilt
\[
    r = \frac{p}{qB}.
\]
Praktisch verwendet man in der Beschleunigerphysik häufig
\[
    p\,[\si{\GeV}/c] \approx 0{,}2998\, B\,[\si{\tesla}]\, r\,[\si{\metre}],
\]
also näherungsweise
\[
    r\,[\si{\metre}] \approx \frac{p\,[\si{\GeV}/c]}{0{,}3\,B\,[\si{\tesla}]}.
\]
\end{formulabox}

\begin{calculationbox}
Für \(\SI{6,9}{\TeV}\)-Protonen ist die kinetische Energie viel größer als die Ruheenergie \(m_pc^2\approx\SI{0,938}{\GeV}\). Daher gilt näherungsweise
\[
    pc \approx E \approx \SI{6,9}{\TeV},
    \qquad
    p\approx \SI{6900}{\GeV}/c.
\]
Mit \(B=\SI{7,9}{\tesla}\) folgt
\[
    r \approx \frac{\num{6900}}{0{,}3\cdot 7{,}9}\,\si{\metre}
    \approx \SI{2911}{\metre}.
\]
Ein Magnetsegment der Länge \(L=\SI{20}{\metre}\) lenkt um den Winkel
\[
    \theta \approx \frac{L}{r}
    = \frac{\SI{20}{\metre}}{\SI{2911}{\metre}}
    \approx \SI{6,87e-3}{\radian}
    \approx \SI{0,394}{\degree}
\]
ab.
Für eine volle Kreisbahn braucht man insgesamt \(2\pi\) rad Ablenkung:
\[
    n = \frac{2\pi}{\theta}
    \approx \frac{2\pi}{6{,}87\cdot10^{-3}}
    \approx 914.
\]
\end{calculationbox}

\begin{answerbox}
Die Ablenkung pro Segment beträgt
\[
    \boxed{\theta\approx \SI{6,9e-3}{\radian}\approx\SI{0,39}{\degree}}.
\]
Für eine vollständige Kreisbahn benötigt man ungefähr
\[
    \boxed{n\approx 9{,}1\cdot 10^2\text{ Segmente}}.
\]
\end{answerbox}

\subsection{Teil (d): Synchrotron am Mondäquator}

\begin{calculationbox}
Der Monddurchmesser sei
\[
    d_{\mathrm{Mond}}=\SI{3500}{\kilo\metre}.
\]
Damit ist der Radius der Kreisbahn
\[
    r = \frac{d_{\mathrm{Mond}}}{2}
    = \SI{1750}{\kilo\metre}
    = \SI{1,75e6}{\metre}.
\]
Mit derselben Magnetfeldstärke \(B=\SI{7,9}{\tesla}\) ergibt sich
\[
    p\,[\si{\GeV}/c]
    \approx 0{,}3\,B\,[\si{\tesla}]\,r\,[\si{\metre}]
    = 0{,}3\cdot 7{,}9\cdot 1{,}75\cdot 10^6.
\]
Also
\[
    p \approx \SI{4,15e6}{\GeV}/c.
\]
Da die Protonen extrem relativistisch wären, gilt wieder \(E\approx pc\):
\[
    E \approx \SI{4,15e6}{\GeV}
    = \SI{4,15e3}{\TeV}
    = \SI{4,15}{\peta\electronvolt}.
\]
\end{calculationbox}

\begin{answerbox}
Ein hypothetisches Mond-Synchrotron mit \(\SI{7,9}{\tesla}\)-Dipolmagneten könnte Protonenenergien von ungefähr
\[
    \boxed{E\approx \SI{4,15}{\peta\electronvolt}}
\]
erreichen.
\end{answerbox}

\subsection{Tafelvortrag zu Aufgabe 1}

\begin{presentationbox}
\textbf{Ziel:} Man zeigt, dass Beschleunigerphysik im Wesentlichen aus Energiegewinn, relativistischer Kinematik und magnetischer Kreisbahn besteht.

\begin{enumerate}[label=\arabic*.]
    \item \textbf{Linearbeschleuniger:} Pro Beschleunigungsspalt gewinnt ein Elektron \(eU=\SI{250}{\kilo\electronvolt}\). Bei \(10^5\) Schritten ergibt das \(\SI{25}{\GeV}\).
    \item \textbf{Halbe Lichtgeschwindigkeit:} Für \(v=c/2\) gilt \(\gamma\approx1{,}1547\), also \(E_{\mathrm{kin}}\approx\SI{79}{\kilo\electronvolt}\). Ein Schritt reicht.
    \item \textbf{Driftröhren:} Die Länge ist \(l=\beta c/(2f)\). Für \(\SI{9}{\MeV}\) folgt \(l\approx\SI{9,3}{\centi\metre}\), für \(\SI{9}{\GeV}\) folgt \(l\approx\SI{67,2}{\centi\metre}\).
    \item \textbf{Dipolmagnet:} Aus \(r=p/(qB)\) beziehungsweise \(p\approx0{,}3Br\) folgt bei \(\SI{6,9}{\TeV}\): \(r\approx\SI{2911}{\metre}\), \(\theta\approx\SI{6,9e-3}{\radian}\), also etwa \(914\) Segmente.
    \item \textbf{Mondring:} Setzt man \(r=\SI{1,75e6}{\metre}\), ergibt sich \(E\approx\SI{4,15}{\peta\electronvolt}\).
\end{enumerate}
\end{presentationbox}

\FloatBarrier
\newpage

%%%%%%%%%%%%%%%%%%%%%%%%%%%%%%%%%%%%%%%%%%%%%%%%%%%%%%%%%%%%%%%%%%%
\section{Aufgabe 2: Luminosität des LHC}

\begin{taskbox}
Die Luminosität beschreibt, wie viele Kollisionen ein Beschleuniger pro Zeit und Wirkungsquerschnitt liefern kann. Gegeben ist eine integrierte Luminosität von etwa \(\lum_{\mathrm{int}}\approx\SI{500}{\per\femtobarn}\) bis Ende 2025 und ein Wirkungsquerschnitt \(\sigma=5\cdot10^{-2}\,\si{\barn}\). Außerdem soll eine Grafik der jährlichen integrierten LHC-Luminosität interpretiert werden.
\end{taskbox}

\subsection{Grundidee: Luminosität und Ereigniszahl}

\begin{definitionbox}
Die momentane Luminosität ist
\[
    \lum(t)=N\phi(t),
\]
wobei \(N\) die Zahl der Targetteilchen und \(\phi(t)\) der Teilchenfluss ist. Die integrierte Luminosität ist
\[
    \lum_{\mathrm{int}} = \int \lum(t)\,\dd t.
\]
Die erwartete Zahl der Ereignisse mit Wirkungsquerschnitt \(\sigma\) ist
\[
    N_{\mathrm{Ereignisse}} = \sigma\,\lum_{\mathrm{int}}.
\]
\end{definitionbox}

\subsection{Teil (a): Anzahl der inelastischen Proton-Proton-Reaktionen}

\begin{calculationbox}[title=Einheitenumrechnung]
Ein Femtobarn ist
\[
    \SI{1}{\femtobarn}=10^{-15}\,\si{\barn}.
\]
Daher gilt
\[
    \SI{1}{\barn}=10^{15}\,\si{\femtobarn}.
\]
Der gegebene Wirkungsquerschnitt ist also
\[
    \sigma = 5\cdot10^{-2}\,\si{\barn}
    = 5\cdot10^{-2}\cdot10^{15}\,\si{\femtobarn}
    = 5\cdot10^{13}\,\si{\femtobarn}.
\]
\end{calculationbox}

\begin{calculationbox}[title=Ereigniszahl]
\[
    N_{\mathrm{Ereignisse}}
    = \sigma\,\lum_{\mathrm{int}}
    = \paren{5\cdot10^{13}\,\si{\femtobarn}}
      \paren{500\,\si{\per\femtobarn}}.
\]
\[
    N_{\mathrm{Ereignisse}}
    = 2500\cdot10^{13}
    = 2{,}5\cdot10^{16}.
\]
\end{calculationbox}

\begin{answerbox}
Bis Ende 2025 sind nach dieser Abschätzung ungefähr
\[
    \boxed{N_{\mathrm{Ereignisse}}\approx 2{,}5\cdot10^{16}}
\]
inelastische Proton-Proton-Reaktionen aufgetreten.
\end{answerbox}

\subsection{Teil (b): Interpretation der Luminositätsgrafik}

\begin{solutionbox}[title=Frage 1: In welchem Jahr fanden die meisten \(pp\)-Wechselwirkungen statt?]
Die Zahl der \(pp\)-Wechselwirkungen ist proportional zur integrierten Luminosität. Daher muss man in der Grafik nur vergleichen, welche Jahreskurve den höchsten Endwert erreicht. Die Kurve für 2024 erreicht mit ungefähr \(\SI{90}{\per\femtobarn}\) den größten Wert.

Also fanden im Jahr
\[
    \boxed{2024}
\]
die bisher meisten \(pp\)-Wechselwirkungen statt.
\end{solutionbox}

\begin{solutionbox}[title=Frage 2: Warum steigen die Verläufe nicht kontinuierlich an?]
Die integrierte Luminosität wächst nur dann, wenn tatsächlich Kollisionen stattfinden und Daten genommen werden. In Zeiten ohne Physikbetrieb bleibt sie konstant. Ursachen für flache Abschnitte sind zum Beispiel Wartung, technische Unterbrechungen, Maschinenstudien, erneutes Befüllen des Rings oder Phasen ohne stabile Strahlen.

Die Kurven steigen deshalb stückweise an und besitzen dazwischen Plateaus.
\end{solutionbox}

\begin{solutionbox}[title=Frage 3: Könnte so ein Verlauf abfallen?]
Nein. Die integrierte Luminosität ist eine aufsummierte Größe:
\[
    \lum_{\mathrm{int}}(t)=\int_{t_0}^{t}\lum(t')\,\dd t'.
\]
Da die Luminosität \(\lum(t)\) physikalisch nicht negativ ist, kann die integrierte Luminosität nur wachsen oder konstant bleiben. Sie kann nicht abfallen.
\end{solutionbox}

\subsection{Tafelvortrag zu Aufgabe 2}

\begin{presentationbox}
\begin{enumerate}[label=\arabic*.]
    \item \textbf{Definition:} Luminosität ist der Proportionalitätsfaktor zwischen Wirkungsquerschnitt und Ereignisrate. Für die Gesamtzahl gilt \(N=\sigma\lum_{\mathrm{int}}\).
    \item \textbf{Einheiten:} \(1\,\si{\barn}=10^{15}\,\si{\femtobarn}\). Deshalb ist \(5\cdot10^{-2}\,\si{\barn}=5\cdot10^{13}\,\si{\femtobarn}\).
    \item \textbf{Ereigniszahl:} \(500\,\si{\per\femtobarn}\cdot5\cdot10^{13}\,\si{\femtobarn}=2{,}5\cdot10^{16}\).
    \item \textbf{Grafik:} 2024 hat den höchsten Endwert, also die meisten \(pp\)-Wechselwirkungen.
    \item \textbf{Plateaus:} Keine Kollisionen bedeutet keine neue integrierte Luminosität.
    \item \textbf{Kein Abfallen:} Eine integrierte nichtnegative Größe ist monoton wachsend oder konstant.
\end{enumerate}
\end{presentationbox}

\FloatBarrier
\newpage

%%%%%%%%%%%%%%%%%%%%%%%%%%%%%%%%%%%%%%%%%%%%%%%%%%%%%%%%%%%%%%%%%%%
\section{Lernzusammenfassung}

\begin{summarybox}[title=Zentrale Formeln]
\begin{align*}
    \Delta E_{\mathrm{kin}} &= qU,\\
    E_{\mathrm{kin}} &= (\gamma-1)mc^2,\\
    \gamma &= \frac{1}{\sqrt{1-v^2/c^2}},\\
    \beta &= \sqrt{1-\frac{1}{\gamma^2}},\\
    l_{\mathrm{Drift}} &= \frac{\beta c}{2f},\\
    r &= \frac{p}{qB},\\
    p\,[\si{\GeV}/c] &\approx 0{,}3\,B\,[\si{\tesla}]\,r\,[\si{\metre}],\\
    N_{\mathrm{Ereignisse}} &= \sigma\lum_{\mathrm{int}}.
\end{align*}
\end{summarybox}

\begin{learningbox}[title=Typische Denkfehler]
\begin{itemize}
    \item \textbf{MeV und GeV verwechseln:} \(10^5\cdot\SI{250}{\kilo\electronvolt}=\SI{25}{\GeV}\), nicht \(\SI{25}{\MeV}\).
    \item \textbf{Zu lange klassisch rechnen:} Bei \(v\sim0{,}5c\) sollte man relativistisch rechnen.
    \item \textbf{Driftröhrenlänge ohne halbe Periode:} Für diesen Typ Linac wird hier \(l=vT/2\) verwendet, nicht \(vT\).
    \item \textbf{Integrierte Luminosität als momentane Rate lesen:} Die integrierte Luminosität ist bereits aufsummiert. Plateaus bedeuten keine neuen Daten, nicht negative Daten.
    \item \textbf{Barn und Femtobarn falsch herum umrechnen:} \(1\,\si{\barn}=10^{15}\,\si{\femtobarn}\), weil ein Femtobarn kleiner ist.
\end{itemize}
\end{learningbox}

\begin{summarybox}[title=Endergebnisse]
\begin{center}
\begin{tabularx}{\textwidth}{@{}lX@{}}
\toprule
Teilaufgabe & Ergebnis \\
\midrule
1(a), Energie & \(E_{\mathrm{kin}}\approx\SI{25}{\GeV}\) \\
1(a), \(v\approx c/2\) & nach einem Beschleunigungsschritt \\
1(b)(i) & \(l\approx\SI{0,0929}{\metre}=\SI{9,3}{\centi\metre}\) \\
1(b)(ii) & \(l\approx\SI{0,672}{\metre}=\SI{67,2}{\centi\metre}\) \\
1(c), Ablenkung & \(\theta\approx\SI{6,9e-3}{\radian}\approx\SI{0,39}{\degree}\) \\
1(c), Segmente & \(n\approx914\) \\
1(d) & \(E\approx\SI{4,15}{\peta\electronvolt}\) \\
2(a) & \(N\approx2{,}5\cdot10^{16}\) inelastische \(pp\)-Reaktionen \\
2(b) & 2024; Plateaus durch Pausen; kein Abfallen möglich \\
\bottomrule
\end{tabularx}
\end{center}
\end{summarybox}

\end{document}
